\documentclass[11pt]{article}
\usepackage[english]{babel}
\usepackage{csquotes}
\usepackage[T1]{fontenc}
\usepackage{lmodern}
\usepackage{graphicx}
\usepackage[backend=biber, style=ieee]{biblatex}
\usepackage[hidelinks]{hyperref}

\addbibresource{references.bib} % importuje soubor se seznamem literatury

\title{A Minimal Academic Paper} % title článku, vypíše se pomocí maketitle
\author{Matěj Brnka} % jméno autora, vypíše se pomocí maketitle

\begin{document} % zahahuje dokument

\maketitle % vypíše nadpis, autora a datum

\begin{abstract} % začátek abstraktu
The universe has fascinated humanity for centuries, inspiring
scientific discoveries and fundamental questions about the nature
of space, time, matter, and energy. From the formation of the
earliest galaxies to the mysterious behavior of black holes,
modern astrophysics seeks to explain the physical processes
that govern the evolution of the cosmos. Advances in observational
astronomy, theoretical physics, and computational modeling have
significantly expanded our understanding of the universe,
revealing phenomena that challenge conventional interpretations
of physical reality. % samotný abstrakt
\end{abstract} % ukončení abstraktu


\noindent\textbf{Keywords:} sit, amet, consectetur, adipiscing

% \tableofcontents % vypíše obsah, pro použítí smazat % na začátku řádku
% \newpage % zalomení stránky, pro použítí smazat % na začátku řádku

\section{Introduction}

The universe is a vast and dynamic system containing billions
of galaxies, each composed of stars, planets, and interstellar
matter. One of the most fascinating predictions of Einstein's
theory of general relativity is the existence of black holes,
regions of spacetime where gravity is so strong that even
light cannot escape.

In 2019, the Event Horizon Telescope collaboration presented
the first direct image of a black hole's shadow in the galaxy
Messier 87, providing compelling observational evidence for
the predictions of general relativity
\parencite{eventhorizon2019}.

\printbibliography % výpis použité literatury

\end{document} % ukončení dokumentu